\documentclass[10pt,a4paper]{amsart}
\usepackage[margin=19mm]{geometry}
\usepackage{amsmath,amssymb,mathtools}
\usepackage[scr=rsfs,cal=euler]{mathalfa}
\usepackage{xfrac}
\usepackage[unicode,hidelinks,bookmarks=false]{hyperref}
\newcommand{\ie}{i.e.\ }
\newcommand{\cf}{cf.\ }
\newcommand{\resp}{respectively\ }
\newcommand{\Subsection}{\subsection}
\providecommand{\nobreakdash}{\nobreak\hbox{-}}
\setlength{\emergencystretch}{2em}
\multlinegap0pt
\usepackage{mathtools}
\usepackage{amssymb}
\usepackage{xcolor}
\usepackage{enumerate}
\usepackage[shortlabels]{enumitem}
\usepackage{tikz}
\usepackage{bm}
\usepackage{float}
\newtheorem{lemma}{Lemma}
\title{A structure-preserving implicit-explicit method for a hyperbolic approximation of fourth-order PDEs}
\author{Rahul Barthwal, Rahuldev Ghorai}
\date{Source: arXiv:2609.05939v1 (CC BY 4.0)}
\begin{document}
\maketitle

\noindent\emph{Source and licence.} A structure-preserving implicit-explicit method for a hyperbolic approximation of fourth-order PDEs. Version arXiv:2609.05939v1 (CC BY 4.0), distributed by arXiv under the Creative Commons Attribution 4.0 International licence, http://creativecommons.org/licenses/by/4.0/, as recorded in the arXiv licence metadata for this exact version and shown on the abstract page https://arxiv.org/abs/2609.05939v1. Full licence notice: \texttt{licenses/arxiv-2609.05939-CC-BY-4.0.txt}.\par\smallskip

\noindent\textit{Editorial scope.} Counted unit: the article's lemma lemma:discrete\_inverse (source0.tex lines 726--744). The article's complete original proof of it is delivered with it (source0.tex lines 745--769, one \texttt{proof} environment of 25 lines). No mathematics is written, repaired or re-derived: the statement and the proof are the article's own lines, in the article's own notation, and no retained line is substituted or re-flowed.  No further source-local context is needed: every cross-reference of every retained line resolves inside the two delivered spans.  Nothing was omitted from any retained span, so the package carries no omission record.  No retained line contains a citation, so the wrapper carries no bibliography and no bibliography entry.  No retained line contains a figure environment, an image-inclusion command or drawing code, so no image is delivered and the package declares no asset dependency.  Editorial formatting only, in addition: an AMS \texttt{amsart} preamble that restates the article's own declarations and macros used by the retained lines, namely the environments \texttt{lemma} on separate counters, together with the standard packages those lines need.  The retained environments print in this document as Lemma 1 in order of appearance, whereas the article prints the same items with the numbering of its own full document.  The complete native source of the article as distributed in the arXiv e-print for this exact version is bundled unchanged as \texttt{sources/209476/source0.tex}.
\begin{lemma}[Discrete inverse estimate]
\label{lemma:discrete_inverse}
There exists a constant $C_d>0$, depending only on the space dimension $d$ and on the choice of finite difference operators, such that
\begin{equation}
    \label{eq:scalar_lap-grad}
    {\|\Delta_h f\|}_{2, \Delta}^2
\leq
\frac{C_d}{\Delta x^2}
{\|\nabla_h^+ f\|}_{2, \Delta}^2.
\end{equation}
For the standard Cartesian discretization defined above, one may take $C_d=4d$. The same estimate holds componentwise for vector-valued grid functions:
\begin{equation}
    \label{eq:vector_lap-grad}
    {\|\Delta_h\mathbf v\|}_{2, \Delta}^2
\leq
\frac{C_d}{\Delta x^2}
{\|\nabla_h^+\mathbf v\|}_{2, \Delta}^2.
\end{equation}
\end{lemma}

\begin{proof}
Here we prove this Lemma for $d = 1$ for the sake of simplicity, using
\[
(\Delta_h f)_i=\frac{(\nabla_h^+f)_i-(\nabla_h^+f)_{i-1}}{\Delta x},
\]
and the following standard inequality
\[
(a-b)^2\le 2(a^2+b^2),
\]
we obtain
\begin{align*}
\|\Delta_h f\|_{2,\Delta}^2
&=
\frac{1}{\Delta x^2}\sum_{i=0}^L
\Delta x((\nabla_h^+f)_i-(\nabla_h^+f)_{i-1})^2\\
&\le\frac{2}{\Delta x^2}\sum_{i=0}^{L-1}
\Delta x(\nabla_h^+f)_i^2
+\frac{1}{\Delta x^2}\Delta x (\nabla_h^+f)_{-1}^2+\frac{1}{\Delta x^2}\Delta x(\nabla_h^+f)_L^2.
\end{align*}
Thanks to periodic boundary, we further get
\begin{equation}
\|\Delta_h f\|_{2,\Delta}^2\le\frac{4}{\Delta x^2}\|\nabla_h^+f\|_{2,\Delta}^2.
\end{equation}
In this case, the estimate \eqref{eq:scalar_lap-grad} holds with $C_d=4$ and a similar procedure can be adapted to prove the estimate for $d>1$. Also the estimate \eqref{eq:vector_lap-grad} for the vector function $\mathbf{v}$ is now straightforward as all its components satisfy \eqref{eq:scalar_lap-grad}.
\end{proof}

\end{document}
